% ======================================================================
% Chapter 2 -- The Designer's Toolkit
% Designing Agentic AI Solutions, Chris Kuo
% Part A: why a notation; the solution-graph notation
% ======================================================================

\chapter{The Designer's Toolkit}
\label{ch:toolkit}

\begin{objectives}
After reading this chapter, you will be able to:
\begin{enumerate}[leftmargin=*,itemsep=2pt]
  \item Draw a solution graph in the book's notation, using its nine node types, five edge types, boundaries and trace tags, and check it against the well-formedness rules.
  \item Apply the seven-step design method to a business process, producing the inputs for a solution graph.
  \item Apply the two patterns used in almost every case, retrieve-then-answer and the approval gate, and recognize the three most common anti-patterns.
  \item Score alternative designs on the eight evaluation criteria, weight the criteria from the business context, and test whether the decision survives a change of weights.
  \item Walk a design node by node to find its failure points, and record the final choice in a decision record that states what would change it.
\end{enumerate}
\end{objectives}

% ----------------------------------------------------------------------
\section{Why a Fixed Notation}
\label{sec:why-notation}

After Harbor Industrial Supply funded Team B's quote design (Chapter~\ref{ch:foundations}), three engineers wrote short descriptions of it for the company wiki. The first wrote that ``the agent looks up pricing and checks policy.'' The second wrote that ``pricing comes from the ERP tool and a rules engine applies the discount cap.'' The third wrote that ``the LLM decides whether a manager needs to approve.''

All three were describing the same system, and the third description was wrong: in Team B's design, code decides when a manager is needed, not the model. But nobody could tell from the prose. The words \emph{agent}, \emph{tool} and \emph{checks} meant different things to each writer. A security reviewer reading the wiki could not tell which components could write to the ERP. A finance reviewer could not tell where the margin floor was enforced. The first real design review took two hours, most of it spent agreeing on what the design was before anyone could ask whether it was good.

This is an old problem in engineering, and the old answer is a shared drawing language. Electrical engineers draw circuits with a fixed set of symbols; software architects use views and diagrams with agreed meanings; process analysts use flowchart and BPMN shapes. A fixed notation does three things prose cannot:
\begin{itemize}[itemsep=1pt]
  \item it makes the \emph{kind} of each component visible at a glance, so the question ``is this a model decision or a rule?'' is answered by a color, not a sentence;
  \item it makes \emph{omissions} visible: a write action with no checkpoint or guardrail in front of it stands out;
  \item it lets two designs be \emph{compared} side by side, which is the heart of the trade-off method later in this chapter.
\end{itemize}

The notation in this chapter is deliberately small. You can learn it in an hour and draw it on a whiteboard with one marker. Every solution graph in the remaining eleven chapters uses it, and the templates in Appendix~C reproduce it for slides and documents.

% ----------------------------------------------------------------------
\section{The Solution-Graph Notation}
\label{sec:notation}

\begin{definition}
A \keyterm{solution graph} is a directed graph whose nodes are the components of a solution---people, triggers, agents, code, MCP servers, knowledge stores, systems of record, checkpoints and guardrails---and whose labeled edges show how control, tool calls, data, events and responsibility move between them.
\end{definition}

The notation has four parts: node types, edge types, boundaries, and annotations. Figure~\ref{fig:legend} shows all of them on one page. The notation uses only three shapes: a box for anything that does work, a cylinder for anything that holds data, and a double-bordered box for a person who must decide. Color tells you which kind of box or cylinder you are looking at, and the label tells you the rest.

\subsection{Nodes}

Table~\ref{tab:nodes} lists the nine node types. Six come straight from Chapter~\ref{ch:foundations}: the four building blocks (code, knowledge, agent, MCP server) plus human checkpoints and guardrails. The other three make the edges of the system explicit: what starts the work (\emph{trigger or user}), where authoritative data lives (\emph{system of record}), and the difference between a top-level agent and a \emph{sub-agent} it directs.

\begin{table}[htbp]
\centering
\small
\caption{The nine node types of the solution-graph notation.}
\label{tab:nodes}
\begin{tabularx}{\textwidth}{@{}>{\raggedright\arraybackslash}p{2.3cm}>{\raggedright\arraybackslash}p{3.1cm}>{\raggedright\arraybackslash}X@{}}
\toprule
Node & Shape & Meaning and required annotation \\
\midrule
Trigger or user & Box, white & A person or event that starts or joins the work: a sales rep, a webhook, a nightly schedule. Annotate the channel. \\
Agent & Box, blue & A model that chooses its next action in a loop (Chapter~\ref{ch:foundations}). Annotate the autonomy level, the stop condition and the model tier. \\
Sub-agent & Box, light blue & An agent invoked by another agent, as a tool or by delegation. Same annotations as an agent. \\
Code step & Box, grey & Deterministic logic: a rule, calculation, validation or routing decision. Annotate the rule source (for example, \emph{config v3}). \\
MCP server & Box, green & First line: server name and owning team. Below it: the tools (and, in italics, resources) that \emph{this} solution may use. \\
Knowledge store & Cylinder, amber & Indexed text or records used for retrieval. Annotate the source owner and freshness. \\
System of record & Dashed cylinder, white & The authoritative database or application (ERP, CRM, core banking). The solution never touches it except through an MCP server. \\
Human checkpoint & Box, pink, double border & A person who approves, edits or rejects. Annotate the role and the service-level target (for example, \emph{4 business hours}). \\
Guardrail & Box, orange & An automated check that blocks or repairs inputs or outputs that violate a policy. Annotate what it checks and what happens on failure. \\
\bottomrule
\end{tabularx}
\end{table}

Two choices in the table deserve comment.

The \emph{MCP server is drawn with its tools listed}, because the tool list is where most security and reliability decisions live. A server may expose twenty tools; the solution graph lists only the ones this solution is allowed to use. That list is the allow-list you will pass to ADK's \texttt{tool\_filter} (Chapter~\ref{ch:foundations}, Listing~1.3), so the diagram and the code can be checked against each other.

The \emph{system of record is drawn separately from the MCP server} in front of it. This keeps a hard rule visible: agents and code steps reach company systems only through a server with a defined contract and owner. An edge from an agent straight to a system of record is a notation error and, almost always, a design error.

\subsection{Edges}

Edges carry five kinds of relationship (Table~\ref{tab:edges}). Every edge gets a label unless its meaning is obvious from the two nodes it joins. Conditions go in square brackets, so a reader can tell a branch from a description: \texttt{[discount > 15\%]}, \texttt{[confidence < 0.8]}, \texttt{[rejected]}.

\begin{table}[htbp]
\centering
\small
\caption{The five edge types.}
\label{tab:edges}
\begin{tabularx}{\textwidth}{@{}>{\raggedright\arraybackslash}p{2.3cm}>{\raggedright\arraybackslash}p{3.1cm}>{\raggedright\arraybackslash}X@{}}
\toprule
Edge & Line & Meaning \\
\midrule
Control flow & Solid, filled head & Work passes from one step to the next. Branches carry bracketed conditions. \\
Tool call & Solid, open head, blue & A node invokes a tool on an MCP server or calls a sub-agent as a tool; the result returns to the caller. Label with the tool name if the server lists several. \\
Data & Dashed & A read from or write to a store or system of record, or a retrieval from a knowledge store. Label \emph{read}, \emph{write} or the content. \\
Event & Dotted, heavy, red & An asynchronous message: a queue message, a webhook, a scheduled tick. Label with the event name. \\
Handoff & Double line & Responsibility for the case moves to another agent or person, who now owns it. \\
\bottomrule
\end{tabularx}
\end{table}

The difference between a tool call and a handoff matters more than it first appears. After a tool call, the caller is still in charge and decides what to do with the result. After a handoff, the caller is done; the receiver owns the case, including its errors. Chapter~\ref{ch:case-airline} turns on this distinction.

\subsection{Boundaries and trace tags}

Two kinds of boundary group nodes. An \keyterm{ownership boundary} (grey dashed) surrounds the nodes one team builds and runs; the pricing team owns the pricing MCP server and its system of record, for instance. A \keyterm{trust boundary} (red dash-dot) surrounds content from a source the solution cannot fully trust: e-mail from outside the company, documents uploaded by customers, web pages. Any edge that crosses a trust boundary into an agent is a place to ask what an attacker could put in that content. Chapter~\ref{ch:case-kyc} returns to trust boundaries in depth.

Finally, a \keyterm{trace tag} is a small black label on a node naming the KPI (\textsf{K1}, \textsf{K2}, \ldots) or risk (\textsf{R1}, \textsf{R2}, \ldots) that justifies it. The tags implement a rule from the book's approach: \emph{every node must trace back to a KPI or a risk}. A node with no tag is either missing its justification or should not be there. Chapter~\ref{ch:business} shows how to choose the KPIs; in this chapter we use them as given.

\begin{figure}[htbp]
\centering
\begin{tikzpicture}[x=1cm,y=1cm]
% nodes, three columns
\node[sg trigger] (n1) at (0,0) {Trigger / user};
\node[sg agent] (n2) at (5,0) {Agent\\[-1pt]{\scriptsize L4, stop: answer}};
\node[sg subagent] (n3) at (10,0) {Sub-agent};
\node[sg code] (n4) at (0,-1.9) {Code step\\[-1pt]{\scriptsize rule: config v3}};
\node[sg mcp] (n5) at (5,-1.9) {MCP server (owner)\nodepart{two}{\scriptsize\ttfamily tool\_a, tool\_b}};
\node[sg store] (n6) at (10,-1.9) {Knowledge\\store};
\node[sg sor] (n7) at (0,-4.1) {System of\\record};
\node[sg human] (n8) at (5,-4.1) {Human checkpoint\\[-1pt]{\scriptsize role, SLA}};
\node[sg guard] (n9) at (10,-4.1) {Guardrail};
\sgtag{n2}{K1}
\sgtag{n8}{R2}
% edges legend
\begin{scope}[yshift=-6.1cm]
\draw[sg flow] (-1,0) -- (0.6,0) node[right, font=\footnotesize] {control flow \texttt{[cond]}};
\draw[sg call] (4,0) -- (5.6,0) node[right, font=\footnotesize] {tool call};
\draw[sg data] (8.4,0) -- (10,0) node[right, font=\footnotesize] {data};
\draw[sg event] (-1,-0.8) -- (0.6,-0.8) node[right, font=\footnotesize] {event};
\draw[sg handoff] (4,-0.8) -- (5.6,-0.8) node[right, font=\footnotesize] {handoff};
\node[sg tag] at (8.7,-0.8) {K1};
\node[right, font=\footnotesize] at (9,-0.8) {trace tag};
\end{scope}
% boundaries
\node[sg boundary, fit=(n4)(n7), inner sep=2.5mm] (ob) {};
\node[sg btitle] at (ob.north west) {ownership boundary};
\node[sg trust, fit=(n3)(n6), inner sep=2.5mm] (tb) {};
\node[sg btitle, text=BrickRed] at (tb.north west) {trust boundary};
\end{tikzpicture}
\caption{The solution-graph notation on one page: nine node types, five edge types, two boundary types and trace tags. The boundaries here only illustrate the line styles.}
\label{fig:legend}
\end{figure}

\subsection{Well-formedness rules}

A drawing in the notation is \emph{well formed} when it passes the eight checks below. They are mechanical on purpose: a reviewer, or the companion design coach, can apply them before discussing whether the design is any good. Each one catches a mistake that recurs in student designs and, unfortunately, in production systems.

\begin{enumerate}[label=\textbf{W\arabic*.},leftmargin=2.5em,itemsep=2pt]
  \item \emph{Every agent and sub-agent states a stop condition and an autonomy level.} An agent without a stop condition has no defined end and no bound on cost.
  \item \emph{Every write action is protected.} Any tool that changes a system of record, sends a message outside the company, or moves money has a human checkpoint or a guardrail on the path before it, or a note explaining why neither is needed.
  \item \emph{No agent touches a system of record directly.} All access goes through an MCP server with a listed tool set.
  \item \emph{Rules live in code steps.} Thresholds, limits and eligibility rules appear in code nodes, not in an agent's annotation.
  \item \emph{Every knowledge store names an owner and a freshness target.}
  \item \emph{Every edge into an agent that crosses a trust boundary passes a guardrail} or is justified in a note.
  \item \emph{Every branch condition is written on its edge}, and the branches out of a node cover all cases (include an \texttt{[else]}).
  \item \emph{Every node has at least one trace tag.}
\end{enumerate}

\begin{designnote}
Well-formed is not the same as good. A well-formed graph can still be the wrong design: too costly, too slow, or solving the wrong problem. The rules guarantee only that the design is explicit enough to be judged. The judging is what Sections~\ref{sec:criteria} to~\ref{sec:decision-record} are for.
\end{designnote}

\subsection{The quote design, redrawn}

Figure~\ref{fig:quote-graph} redraws Harbor's quote design from Chapter~\ref{ch:foundations} in the notation, with these KPIs and risks:
\begin{itemize}[itemsep=0pt]
  \item[\textsf{K1}] median quote turnaround below 4 hours (today 26);
  \item[\textsf{K2}] average realized margin on quoted orders held at or above today's level;
  \item[\textsf{K3}] pricing-analyst minutes per quote below 8 (today about 35);
  \item[\textsf{R1}] a quote issued below the margin floor;
  \item[\textsf{R2}] the wrong customer or contract applied to a quote;
  \item[\textsf{R3}] a number in the cover note that disagrees with the quote itself.
\end{itemize}

\begin{figure}[htbp]
\centering
\begin{tikzpicture}[x=0.97cm,y=1cm, sg base/.append style={text width=2.05cm}]
% row 1
\node[sg trigger] (rep) at (0,0) {Sales rep\\[-1pt]{\scriptsize e-mail / chat}};
\node[sg guard] (g1) at (3.1,0) {Input\\screen};
\node[sg agent] (intake) at (6.2,0) {Intake agent\\[-1pt]{\scriptsize L1; stop: schema filled}};
\node[sg code] (fetch) at (9.3,0) {Validate \&\\fetch prices};
\node[sg mcp, text width=2.75cm] (erp) at (12.45,0) {ERP server\\[-1pt]{\scriptsize owner: Pricing}\nodepart{two}{\scriptsize\ttfamily get\_item\_pricing\\get\_contract\_terms}};
% row 2
\node[sg store, text width=1.9cm] (kb) at (3.1,-2.7) {Policy \& contract index\\[-1pt]{\scriptsize Finance; daily}};
\node[sg code] (policy) at (9.3,-2.7) {Policy engine\\[-1pt]{\scriptsize cap 15\%, floor 22\%\\config v3}};
\node[sg sor] (erpdb) at (12.45,-2.7) {ERP};
% row 3
\node[sg human] (mgr) at (5.4,-5.2) {Sales manager\\[-1pt]{\scriptsize SLA 4 business h}};
\node[sg agent] (draft) at (9.3,-5.2) {Drafting agent\\[-1pt]{\scriptsize L1; stop: note written}};
% row 4
\node[sg trigger] (rep2) at (0,-7.6) {Sales rep\\[-1pt]{\scriptsize reviews, sends}};
\node[sg mcp, text width=2.75cm] (crm) at (5.4,-7.6) {CRM server\\[-1pt]{\scriptsize owner: Sales Ops}\nodepart{two}{\scriptsize\ttfamily create\_quote\_draft}};
\node[sg guard] (g2) at (9.3,-7.6) {Numbers\\match?};
\node[sg sor] (crmdb) at (12.45,-7.6) {CRM};
% edges
\draw[sg flow] (rep) -- (g1);
\draw[sg flow] (g1) -- (intake);
\draw[sg flow] (intake) -- (fetch);
\draw[sg call] (fetch) -- (erp);
\draw[sg data] (erp) -- (erpdb);
\draw[sg flow] (fetch) -- node[sg lab, right] {[complete]} (policy);
\draw[sg flow] (fetch.north) to[bend right=20] node[sg lab, above] {[missing fields]: ask rep} (rep.north);
\draw[sg data] (kb) -- node[sg lab, above] {read: policy \S, version} (policy);
\draw[sg flow] (policy) -- node[sg lab, sloped, above] {[exception]} (mgr);
\draw[sg flow] (policy) -- node[sg lab, right] {[within policy]} (draft);
\draw[sg flow] (mgr) -- node[sg lab, above] {[approved]} (draft);
\draw[sg flow] (mgr.west) -| node[sg lab, pos=0.75, left] {[rejected]} (rep2.north);
\draw[sg flow] (draft) -- (g2);
\draw[sg flow] (g2) -- node[sg lab, above] {[pass]} (crm);
\draw[sg data] (crm.south) to[bend right=25] node[sg lab, below] {write: quote draft} (crmdb.south);
\draw[sg handoff] (crm) -- (rep2);
% boundaries
\node[sg trust, fit=(rep)] (tb) {};
\node[sg btitle, text=BrickRed] at (tb.north west) {untrusted text};
\node[sg boundary, fit=(erp)(erpdb)] (pb) {};
\node[sg btitle] at (pb.north west) {Pricing team};
% tags
\sgtag{intake}{K1 K3}
\sgtag{g1}{R2}
\sgtag{fetch}{R2}
\sgtag{policy}{K2 R1}
\sgtag{mgr}{R1}
\sgtag{draft}{K3}
\sgtag{g2}{R3}
\sgtag{kb}{K2}
\sgtag{crm}{K1}
\end{tikzpicture}
\caption{Harbor's quote design as a solution graph. Compare Figure~1.1: the same design, but every component's kind, owner, rule source and justification is now on the page.}
\label{fig:quote-graph}
\end{figure}

Compare it with the informal sketch in Figure~1.1. Three things are now explicit that were not before. The model never decides whether a manager is needed: that branch leaves a code node. The agents' autonomy is bounded: both are Level~1 with a defined stop. And the one action that changes a system of record, writing the draft to the CRM, sits behind a guardrail that checks every number in the cover note against the policy engine's output, followed by a handoff to the rep, who sends.

Run the eight checks against the figure as an exercise. You should find that it passes all of them except W8 for the two systems of record and the rep nodes; by convention, triggers, users and systems of record are exempt from W8, because they exist whether or not the solution does.

\subsection{Three levels of detail}

One diagram cannot hold everything. We use three levels, each zooming into the one above:
\begin{description}[leftmargin=1em,itemsep=2pt]
  \item[Context view.] One box for the solution, surrounded by the people, triggers and systems it touches. Used in the business brief, before any design choice is made.
  \item[Solution graph.] The level of Figure~\ref{fig:quote-graph}. Every alternative design in every case is drawn at this level, so alternatives can be compared.
  \item[Node card.] A short specification for one node, in a fixed format (Table~\ref{tab:nodecard}). Agents, MCP servers and guardrails always get a card; simple code steps may not need one.
\end{description}

\begin{table}[htbp]
\centering
\small
\caption{A node card for the intake agent in Figure~\ref{fig:quote-graph}.}
\label{tab:nodecard}
\begin{tabularx}{\textwidth}{@{}>{\raggedright\arraybackslash}p{3cm}>{\raggedright\arraybackslash}X@{}}
\toprule
Field & Intake agent \\
\midrule
Purpose & Turn a rep's free-text request into a structured quote request. \\
Inputs & Message text and sender ID, after the input screen. \\
Outputs & JSON matching \texttt{QuoteRequest} schema: customer ID, line items (SKU, quantity), requested discount, needed-by date; or a list of missing fields. \\
Autonomy and stop & Level 1; one model call with structured output; stops when the schema is filled or the missing-fields list is produced. \\
Tools & None. Matching the customer name to a customer ID is left to the validate-and-fetch code step, which can be tested exactly. \\
Model tier & Small, fast model; structured-output support required. \\
Quality target & 98\% field-level accuracy on the 300-message evaluation set; zero invented SKUs. \\
Cost and latency budget & Below \$0.01 and 3 seconds per request at the 95th percentile. \\
Failure handling & Schema failure: retry once, then route to a pricing analyst. Customer name unclear: return the name as written; the code step asks the rep rather than guessing. \\
Traces to & \textsf{K1}, \textsf{K3}; mitigates \textsf{R2} with the customer check. \\
\bottomrule
\end{tabularx}
\end{table}

The node card's \emph{quality target} and \emph{cost and latency budget} rows are what make a design testable. Chapter~\ref{ch:case-fraud} shows how to build the evaluation set that the quality target refers to.
